% Diagrama TikZ: Arquitectura jerárquica 3 niveles — v4 legibilidad IEEE
% Mejoras v4 (REORG FASE C, resuelve concern C6):
%   - Fuentes elevadas: \tiny -> \footnotesize; \scriptsize -> \small
%   - Rellenos de bandas saturados: !3 / !4  -> !12 / !15
%   - Rellenos de cajas de métricas saturados: !4 -> !22
%   - text width de métricas ampliado: 2.6cm -> 4.2cm (evita word-wrap)
%   - Canvas ampliado a 22cm para que \resizebox no aplaste las fuentes
%   - Conservadas todas las flechas/conexiones y labels (mismo \label)

\begin{figure}[H]
\centering
\resizebox{0.97\textwidth}{!}{%
\begin{tikzpicture}[
  % ── Paleta por nivel ──
  clrN1/.style={green!55!black},
  clrN2/.style={blue!55!black},
  clrN3/.style={orange!55!black},
  % ── Bandas de nivel ──
  levelband/.style={rounded corners=6pt, line width=1.2pt, inner sep=10pt},
  bandN1/.style={levelband, draw=green!50!black, fill=green!12,
                 drop shadow={shadow xshift=0.5pt, shadow yshift=-0.5pt, fill=black!8}},
  bandN2/.style={levelband, draw=blue!45!black, fill=blue!10,
                 drop shadow={shadow xshift=0.5pt, shadow yshift=-0.5pt, fill=black!8}},
  bandN3/.style={levelband, draw=orange!45!black, fill=orange!15,
                 drop shadow={shadow xshift=0.5pt, shadow yshift=-0.5pt, fill=black!8}},
  % ── Etiqueta de nivel ──
  lvllabel/.style={font=\normalsize\bfseries\sffamily, anchor=north west},
  % ── Componentes (base) ──
  comp/.style={rounded corners=3pt, minimum height=0.95cm,
               font=\footnotesize\sffamily, inner sep=4pt, align=center,
               line width=0.7pt,
               drop shadow={shadow xshift=0.4pt, shadow yshift=-0.4pt, fill=black!10}},
  % Variantes de color (rellenos saturados para impresión)
  meter/.style=  {comp, draw=brown!55, fill=brown!22,  text=brown!70!black},
  znode/.style=  {comp, draw=green!55!black, fill=green!22, text=green!40!black},
  otbrst/.style= {comp, draw=green!60!black, fill=green!28, text=green!35!black},
  halowst/.style={comp, draw=orange!70!black, fill=orange!25, text=orange!40!black},
  tbcomp/.style= {comp, draw=blue!55, fill=blue!20, text=blue!40!black},
  dbcomp/.style= {comp, draw=purple!55, fill=purple!20, text=purple!40!black},
  vizcomp/.style={comp, draw=teal!55, fill=teal!22, text=teal!40!black},
  srvcomp/.style={comp, draw=red!45!black, fill=red!18, text=red!45!black},
  dashcomp/.style={comp, draw=gray!55, fill=gray!20, text=gray!50!black},
  % ── Flechas ──
  arr/.style={-{Latex[length=2.5mm, width=1.8mm]}, line width=0.8pt},
  arrbold/.style={-{Latex[length=3mm, width=2.2mm]}, line width=1.3pt},
  % ── Anotación sobre flecha ──
  ann/.style={font=\footnotesize\sffamily, text=#1},
  % ── Caja de métricas (más ancha y con relleno saturado) ──
  metricbox/.style={draw=#1!50, fill=#1!22, rounded corners=3pt,
                    font=\footnotesize\sffamily, inner sep=5pt, text width=4.2cm,
                    anchor=west, text=#1!60!black, line width=0.6pt},
  % ── Grupo funcional ──
  funcgroup/.style={draw=#1!40, fill=#1!8, rounded corners=4pt,
                    line width=0.6pt, inner sep=6pt},
]

% ═════════════════════════════════════════════════════════════════
%  NIVEL 1 — Red de Campo (malla Thread)
% ═════════════════════════════════════════════════════════════════
\node[bandN1, minimum width=21cm, minimum height=3cm,
      anchor=south west] (L1) at (0, 0) {};
\node[lvllabel, clrN1] at (0.35, 2.85)
  {Nivel 1 \normalfont\sffamily\small --- Red de campo (malla Thread)};

% Medidores y nodos (3 explícitos + dots + nodo N)
\foreach \i/\x in {1/1.8, 2/5.0, 3/8.2} {
  \node[meter, minimum width=1.7cm] (M\i) at (\x, 0.50)
    {\raisebox{1pt}{\scalebox{0.8}{$\mathbf{\sim}$}} Medidor\\[-1pt]DLMS/COSEM};
  \node[znode, minimum width=2.0cm] (N\i) at (\x, 1.70)
    {Nodo Zephyr\\[-1pt]802.15.4};
  \draw[arr, brown!60] (M\i.north) -- (N\i.south)
    node[midway, right=2pt, ann=brown!60!black] {RS-485};
}

% Puntos suspensivos
\node[font=\Large, gray!50] at (10.6, 1.1) {\textbf{$\cdots$}};

% Último nodo
\node[meter, minimum width=1.7cm] (MN) at (12.6, 0.50)
  {\raisebox{1pt}{\scalebox{0.8}{$\mathbf{\sim}$}} Medidor\\[-1pt]DLMS/COSEM};
\node[znode, minimum width=2.0cm] (NN) at (12.6, 1.70)
  {Nodo $n{=}60$\\[-1pt]802.15.4};
\draw[arr, brown!60] (MN.north) -- (NN.south)
  node[midway, right=2pt, ann=brown!60!black] {RS-485};

% Enlaces mesh Thread
\draw[green!45!black, line width=0.9pt, densely dashed,
      decorate, decoration={zigzag, amplitude=0.8pt, segment length=6pt}]
  (N1.east) -- (N2.west);
\draw[green!45!black, line width=0.9pt, densely dashed,
      decorate, decoration={zigzag, amplitude=0.8pt, segment length=6pt}]
  (N2.east) -- (N3.west);
\draw[green!45!black, line width=0.9pt, densely dashed,
      decorate, decoration={zigzag, amplitude=0.8pt, segment length=6pt}]
  (N3.east) -- (10.1, 1.70);
\draw[green!45!black, line width=0.9pt, densely dashed,
      decorate, decoration={zigzag, amplitude=0.8pt, segment length=6pt}]
  (11.1, 1.70) -- (NN.west);

% Cajas de métricas Nivel 1 (text width ampliado a 4.2cm)
\node[metricbox=green] at (15.4, 1.95) {
  \textbf{Pila de protocolos}\\[1pt]
  IEEE 802.15.4 $\rightarrow$ 6LoWPAN\\
  Thread 1.4 $\rightarrow$ CoAP\\
  LwM2M (IPSO Objects)
};
\node[metricbox=green] at (15.4, 0.55) {
  \textbf{1--60 nodos en malla}\\[1pt]
  Latencia: 6\,ms (P50)\\
  Consumo: 27\,mA activo\\
  Autorreparación multisalto
};

% ═════════════════════════════════════════════════════════════════
%  NIVEL 2 — Pasarela de Borde (RPi 4 + Docker)
% ═════════════════════════════════════════════════════════════════
\node[bandN2, minimum width=21cm, minimum height=4cm,
      anchor=south west] (L2) at (0, 3.6) {};
\node[lvllabel, clrN2] at (0.35, 7.45)
  {Nivel 2 \normalfont\sffamily\small --- Pasarela de borde (Raspberry Pi 4 + Docker)};

% OTBR (border router) — posición izquierda
\node[otbrst, minimum width=2.3cm, minimum height=1.0cm]
  (OTBR) at (2.0, 4.5) {OTBR\\[-1pt]nRF52840};

% HaLow STA — posición izquierda arriba
\node[halowst, minimum width=2.3cm, minimum height=1.0cm]
  (HSTA) at (2.0, 6.6) {HaLow STA\\[-1pt]MM6108 SDIO};

% Grupo ThingsBoard Edge
\node[funcgroup=blue, minimum width=7.8cm, minimum height=3cm,
      anchor=south west] (tbgroup) at (4.5, 3.9) {};
\node[font=\small\bfseries\sffamily, blue!55!black,
      anchor=north west] at (4.7, 6.85)
  {ThingsBoard Edge 3.6};

% Componentes TB — fila inferior
\node[tbcomp, minimum width=2.4cm] (LWM2M) at (6.0, 4.55)
  {Transporte\\[-1pt]LwM2M / CoAP};
\node[tbcomp, minimum width=2.4cm] (RULES) at (9.2, 4.55)
  {Motor de reglas\\[-1pt](Rule Engine)};
% Componentes TB — fila superior
\node[tbcomp, minimum width=2.4cm] (DEVMGR) at (6.0, 5.85)
  {Perfiles y\\[-1pt]dispositivos};
\node[tbcomp, minimum width=2.4cm] (SYNC) at (9.2, 5.85)
  {Sincr.\ nube\\[-1pt](gRPC/MQTT)};

% Flechas internas TB (horizontales y verticales cortas — sin cruces)
\draw[arr, blue!35, line width=0.5pt] (LWM2M.east) -- (RULES.west);
\draw[arr, blue!35, line width=0.5pt] (LWM2M.north) -- (DEVMGR.south);
\draw[arr, blue!35, line width=0.5pt] (DEVMGR.east) -- (SYNC.west);
\draw[arr, blue!35, line width=0.5pt] (RULES.north) -- (SYNC.south);

% OTBR → LwM2M (horizontal directa)
\draw[arr, blue!60, line width=0.9pt] (OTBR.east) -- (LWM2M.west)
  node[midway, above=3pt, ann=blue!50!black,
       fill=white, inner sep=2pt, rounded corners=1pt] {CoAP/UDP};

% SYNC → HaLow STA — ruta por ENCIMA del grupo TB (sin cruzar texto)
\draw[arr, teal!60, line width=0.7pt]
  (SYNC.north) -- ++(0, 0.95) -| (HSTA.north)
  node[pos=0.25, above=2pt, ann=teal!50!black,
       fill=white, inner sep=2pt, rounded corners=1pt] {MQTT/TLS};

% PostgreSQL + TimescaleDB
\node[dbcomp, minimum width=2.5cm, minimum height=1.0cm]
  (PG) at (13.7, 4.55) {PostgreSQL\\[-1pt]TimescaleDB};

% Grafana
\node[vizcomp, minimum width=2.5cm, minimum height=1.0cm]
  (GRAF) at (13.7, 5.95) {Grafana\\[-1pt]Tableros};

% Rule → PG (horizontal directa)
\draw[arr, purple!55, line width=0.6pt] (RULES.east) -- (PG.west)
  node[midway, above=3pt, ann=purple!50!black,
       fill=white, inner sep=2pt, rounded corners=1pt] {SQL};
% PG → Grafana (vertical directa)
\draw[arr, teal!50, line width=0.6pt] (PG.north) -- (GRAF.south);

% Cajas de métricas Nivel 2 (text width 4.2cm)
\node[metricbox=blue] at (15.9, 6.25) {
  \textbf{Procesamiento local}\\[1pt]
  Docker: 7 contenedores\\
  RPi 4 (4\,GB) + OpenWrt\\
  Respaldo sin conexión: 7 días
};
\node[metricbox=blue] at (15.9, 4.55) {
  \textbf{Métricas de borde}\\[1pt]
  Latencia E2E: 8\,ms P90\\
  Reducción WAN: 72\,\%
};

% ═════════════════════════════════════════════════════════════════
%  NIVEL 3 — Enlace Troncal HaLow + Servidor
% ═════════════════════════════════════════════════════════════════
\node[bandN3, minimum width=21cm, minimum height=2.8cm,
      anchor=south west] (L3) at (0, 8.2) {};
\node[lvllabel, clrN3] at (0.35, 10.85)
  {Nivel 3 \normalfont\sffamily\small --- Enlace troncal Wi-Fi HaLow + Servidor};

% HaLow Mesh (AHPI6108E / Tube AHM)
\node[halowst, minimum width=2.8cm, minimum height=1.05cm]
  (HAP) at (3.4, 9.2) {HaLow 802.11ah\\[-1pt]AHPI6108E / Tube AHM};

% Servidor ThingsBoard
\node[srvcomp, minimum width=3.0cm, minimum height=1.05cm]
  (SRV) at (8.8, 9.2) {ThingsBoard\\[-1pt]Servidor / Nube};

% Panel operador
\node[dashcomp, minimum width=2.8cm, minimum height=1.05cm]
  (DASH) at (13.7, 9.2) {Panel del\\[-1pt]operador};

% Flechas Nivel 3 (horizontales directas)
\draw[arr, orange!65, line width=0.9pt] (HAP.east) -- (SRV.west)
  node[midway, above=3pt, ann=orange!50!black,
       fill=white, inner sep=2pt, rounded corners=1pt] {Ethernet/IP};
\draw[arr, gray!50, line width=0.6pt] (SRV.east) -- (DASH.west)
  node[midway, above=3pt, ann=gray!50!black,
       fill=white, inner sep=2pt, rounded corners=1pt] {HTTPS / WS};

% Cajas de métricas Nivel 3 (text width 4.2cm)
\node[metricbox=orange] at (15.9, 9.85) {
  \textbf{Wi-Fi HaLow (802.11ah)}\\[1pt]
  Sub-1\,GHz: 902--928\,MHz\\
  Canal 14, WPA3-SAE
};
\node[metricbox=orange] at (15.9, 8.75) {
  \textbf{Enlace troncal @ 2\,MHz}\\[1pt]
  DL: 3{,}77\,Mbps TCP\\
  UL: 1{,}34\,Mbps TCP\\
  Latencia: 8{,}5\,ms med.
};

% ═════════════════════════════════════════════════════════════════
%  CONEXIONES ENTRE NIVELES — rutas sin cruces con textos
% ═════════════════════════════════════════════════════════════════

% Nivel 1 → Nivel 2 (Thread → OTBR) — vertical directa
\draw[arrbold, green!55!black]
  ([yshift=2pt] N3.north -| OTBR) -- (OTBR.south)
  node[pos=0.45, right=3pt,
       font=\small\bfseries\sffamily, green!40!black,
       fill=white, inner sep=2pt, rounded corners=1pt] {6LoWPAN};

% Nivel 2 → Nivel 3 (HaLow STA → HaLow AP) — vertical directa
\draw[arrbold, orange!60!black]
  (HSTA.north) -- ([yshift=-2pt]HAP.south)
  node[pos=0.45, right=3pt,
       font=\small\bfseries\sffamily, orange!45!black,
       fill=white, inner sep=2pt, rounded corners=1pt] {802.11ah};

\end{tikzpicture}%
}% end resizebox

\caption[Arquitectura jerárquica de tres niveles]{%
Arquitectura jerárquica de tres niveles propuesta para infraestructura
de medición avanzada (AMI) con procesamiento de borde.
\textbf{Nivel~1:} 1--60 nodos Zephyr con radio IEEE~802.15.4 en red
Thread (6LoWPAN/CoAP/LwM2M) conectados a medidores vía RS-485
DLMS/COSEM.
\textbf{Nivel~2:} Pasarela Raspberry Pi~4 ejecutando ThingsBoard
Edge~3.6 con transporte LwM2M, motor de reglas local,
PostgreSQL+TimescaleDB y Grafana; enlace HaLow STA (IEEE~802.11ah)
hacia Nivel~3.
\textbf{Nivel~3:} Enlace troncal Wi-Fi HaLow (sub-1\,GHz, 909\,MHz)
vía Morse Micro AP hacia servidor ThingsBoard.}
\label{fig:arquitectura-tres-niveles}
\end{figure}
